\documentclass[11pt]{article}
\usepackage[T1]{fontenc}
\usepackage{lmodern,microtype,amsmath,amssymb,amsthm,booktabs,longtable,array}
\usepackage[margin=1in]{geometry}
\usepackage{xurl}
\usepackage[hidelinks]{hyperref}
\setlength{\emergencystretch}{3em}
\newtheorem{theorem}{Theorem}
\newtheorem{lemma}[theorem]{Lemma}
\newtheorem{proposition}[theorem]{Proposition}
\providecommand{\tightlist}{\setlength{\itemsep}{0pt}\setlength{\parskip}{0pt}}
\providecommand{\passthrough}[1]{#1}
\hypersetup{pdfauthor={Yongxian Zhang}}

\title{An explicit positive-definite counterexample to half-line q-permanent monotonicity}
\author{Yongxian Zhang\\[2pt]\small School of Computer Science and Engineering\\\small South China University of Technology\\\small\texttt{mxymmxym1@gmail.com}\\\small ORCID: 0009-0000-3864-3536}
\date{8 October 2026}
\begin{document}\maketitle
\begin{abstract}We exhibit a rational real symmetric positive-definite matrix of order four whose inversion-weighted q-permanent decreases between q=49 and q=50. Exact principal minors establish positive definiteness, and a six-term permutation computation gives the complete polynomial. The counterexample disproves da Fonseca's positive-definite half-line extension, independently of the convention at its left endpoint. The actual matrix and q-permanent counterexample are completely formalized in Lean.\end{abstract}
\section{The half-line conjecture}\label{the-half-line-conjecture}

For a real matrix \(A=(a_{ij})\) of order \(n\), define \[
 P_q(A)=\sum_{\sigma\in S_n}q^{\operatorname{inv}(\sigma)}
                  \prod_{i=1}^n a_{i,\sigma(i)},\qquad
 \operatorname{inv}(\sigma)=\#\{(i,j):i<j,\ \sigma(i)>\sigma(j)\}.
\] The index ordering is part of the definition. Da Fonseca's half-line
extension asks whether, for a real symmetric positive-definite
non-diagonal matrix, there is a left endpoint at or below \(-1\) from
which this function is strictly increasing up to \(+\infty\). His 2018
formulation uses an endpoint strictly below \(-1\). A decrease between
49 and 50 refutes both versions, independently of the left endpoint
convention.

\section{The exact counterexample}\label{the-exact-counterexample}

\textbf{Theorem.} The rational matrix \[
A=\begin{pmatrix}
1&-99/100&0&1/500\\
-99/100&1&0&1/8\\
0&0&1&0\\
1/500&1/8&0&1
\end{pmatrix}
\] is positive definite, and \(P_{49}(A)>P_{50}(A)\). Consequently it
disproves the positive-definite half-line conjecture.

\textbf{Proof.} Its successive leading principal minors are \[
1,\quad\frac{199}{10000},\quad\frac{199}{10000},\quad\frac{59}{15625}.
\] They are strictly positive, so Sylvester's criterion proves positive
definiteness. A permutation can contribute only if it fixes index 3. The
six permutations of indices \(\{1,2,4\}\) give the following complete
list.

\begin{longtable}[]{@{}lrl@{}}
\toprule\noalign{}
Permutation & Inversions & Monomial contribution \\
\midrule\noalign{}
\endhead
\bottomrule\noalign{}
\endlastfoot
Identity & 0 & \(1\) \\
\((12)\) & 1 & \((9801/10000)q\) \\
\((24)\) & 3 & \((1/64)q^3\) \\
\((14)\) & 5 & \((1/250000)q^5\) \\
\((124)\) & 4 & \(-(99/400000)q^4\) \\
\((142)\) & 4 & \(-(99/400000)q^4\) \\
\end{longtable}

Thus exact enumeration, retaining all four index positions, gives \[
P_q(A)=1+\frac{9801}{10000}q+\frac1{64}q^3
                -\frac{99}{200000}q^4+\frac1{250000}q^5.
\] Evaluation yields \[
P_{49}(A)=\frac{81807513}{500000},\qquad
P_{50}(A)=\frac{7969}{50},\qquad
P_{49}(A)-P_{50}(A)=\frac{2117513}{500000}>0.
\] Since \(49<50\), the theorem follows. No numerical optimization,
approximate eigenvalue or derivative test is used. \(\square\)

\section{Scope and reproduction}\label{scope-and-reproduction}

This is the positive-definite half-line conjecture, not Bapat's original
monotonicity conjecture restricted to \([-1,1]\). Deleting the isolated
identity index or moving the nonzero triangle to three consecutive
indices would change the inversion statistic and the displayed
polynomial. The checker \texttt{verify\_exact.py} independently
enumerates all permutations and evaluates every determinant and
polynomial value in rational arithmetic. The accompanying Lean proof
works with the actual matrix positive-definiteness definition and the
actual inversion-weighted permutation sum, rather than assuming a
surrogate polynomial identity. The separate minimum-dimension and
arbitrary-\(t>1\) extensions in the earlier research note are outside
this preprint and its formalization claims.

\section{References}\label{references}

\begin{enumerate}
\def\labelenumi{\arabic{enumi}.}
\tightlist
\item
  C. M. da Fonseca, \emph{The mu-permanent of a tridiagonal matrix,
  orthogonal polynomials, and chain sequences}, Linear Algebra and its
  Applications \textbf{432} (2010), 1258--1266, Conjecture 1, p.~1259.
  \url{https://doi.org/10.1016/j.laa.2009.10.036}.
\item
  C. M. da Fonseca, \emph{The mu-permanent revisited},
  arXiv:1804.02231v1 (2018), Conjecture 2.
  \url{https://arxiv.org/abs/1804.02231v1}.
\item
  L. Mitchell, \emph{A note on Bapat's q-permanent conjecture},
  Operators and Matrices \textbf{14} (2020), 915--919, especially the
  remark on pp.~917--918 distinguishing the positive-definite conjecture
  from its singular extension.
  \url{https://doi.org/10.7153/oam-2020-14-56}.
\end{enumerate}

\section{Formal verification and reproduction}
The principal theorem has a complete Lean proof. The precise scope is:
Only the displayed 4-by-4 real rational positive-definite counterexample and its exact decrease between q=49 and q=50. Minimal dimension and the all-t>1 family from the older research note are outside this paper.
The public proof source snapshot is fixed at commit
\nolinkurl{02af7db99db789fecb0d5f5f8f37b1a4e40762a0} in \url{https://github.com/mxym/math}.
The proof packages, with their own dependency locks, audit records and
reproduction instructions, are:
\begin{itemize}
\item \nolinkurl{formalizations/q-permanent-halfline-counterexample}.
\end{itemize}
The main formal entry points are:
\begin{itemize}
\item \nolinkurl{QPermanentHalfline.explicit_counterexample}.
\item \nolinkurl{QPermanentHalfline.halfLineConjecture_false}.
\end{itemize}
Lean is fixed to version 4.34.1, compiler commit
\nolinkurl{5045d0056413266e57c625dcd7c365b10e377c52}; Mathlib is fixed to
\nolinkurl{d13f23b723b8a846827a245b89c10fc7d3f11612}.
The audited mathematical roots depend only on the standard axioms
\texttt{propext}, \texttt{Classical.choice}, and \texttt{Quot.sound}.
The archived audits state their actual compilation and replay scopes;
publication of this typeset version does not imply a new execution of every
previous empty-kernel replay. The companion README maps each printed
principal result to its formal theorem and explains how to repeat verification.

\section*{Research and preparation disclosures}
This research, formalization and manuscript preparation used AI assistance.
The author is responsible for the manuscript. The principal claims are
supported by the readable proofs and the public Lean artifacts specified
above. Formal verification is not external human peer review; no such review
is claimed. No external funding was received. No worldwide priority or
journal acceptance is asserted.
\end{document}
